\documentclass[track=junior,level=3,teacher]{kidscode}
\begin{document}
\kctitle{Answer Key}{What each activity should look like, and what to look for}

\begin{teachernote}[How to use the answer key]
Answers describe what you should \emph{see} and \emph{hear}. Many projects are correct in more
than one way: accept anything that does what the task asks. For open tasks, use the
\enquote{what to look for} notes.
\end{teachernote}
\kcweeklytoc
\kcinput{answers/week01}
\kcinput{answers/week02}
\kcinput{answers/week03}
\kcinput{answers/week04}
\kcinput{answers/week05}
\kcinput{answers/week06}
\kcinput{answers/week07}
\kcinput{answers/week08}
\kcinput{answers/week09}
\kcinput{answers/week10}
\kcinput{answers/assessments}
\end{document}
